\section{The \texorpdfstring{$\nim{2}$}{*2} Question: Details and Ring Listings}\label{S-star2}
\subsection{The Textbook Candidate}
Fig.~\ref{fig:star2_textbook} shows the candidate \pos{W a2 c2 / B b7 d7} on files a--d. Under the count rule
with the textbook test, it is the only $\nim{k}$ with $k\ge2$ in the census of Section~\ref{sec:atlas}. If white
moves first, a2--a4 leaves $\star$ and c2--c4 leaves $0$; a2--a3 and c2--c3 are worse. The position is symmetric
(turn the board around and swap the colors), so black's d7--d5 leaves $\star$ and b7--b5 leaves $0$. Thus
$G=\game{0,\star}{0,\star}=\nim{2}$.

The weak point is black's plan b7--b5 and b5--b4. After b5--b4 the a-pawn is in the situation of
Fig.~\ref{fig:epS}: the textbook test scores a2--a4 as a passed pawn, $+\BIG$, but in chess black takes
b4$\times$a3 en passant. Under the e.p.-aware test, \pos{W a2 c2 / B b4 d7} is worth $-3$, a black win, so
b7--b5 is worth $-2\pm1$ and, by symmetry, c2--c4 is worth $2\pm1$. The other moves, among them a2--a4 and
d7--d5, which still lead to $\star$, are dominated or reversible. Hence
\[
G=\game{2\pm1}{-2\pm1},
\]
a hot game with mean $0$ and temperature $2$ \lab{C2}.

\subsection{Census and Seven Pawns}
Our CGT engine and an independent C solver, which tests whether $G+\nim{k}$ is a second-player win, agree on
every record, and the independent whole-board solver finds no $\nim{k}$ with $2\le k\le 292$ among
$793{,}299$ labeled placements on widths 1--4 (none reduced by symmetry or translation). Every position in scope has birthday (Section~\ref{sec:notation}) at most $292$, so no
larger nimber can occur. Under the textbook test the census has $6{,}470$ zeros, $3{,}891$ stars and the
candidate; the e.p.-aware test removes that $\nim{2}$. In non-default settings a few other census
positions are $\nim{2}$; Section~\ref{sec:real-cand} tests one of them in real chess.

Both fit beside T, and their pawns can arise in a game: A by a2--a4, b2--b3, c2$\times$d3 and c7--c5, B by
b2--b4, b7$\times$a6, c7--c6 and d7--d5, with g2--g4 and g7--g5 for T; pieces take the other men.

In position A (Fig.~\ref{fig:star2_seven_a}), if white moves first, a4--a5 leaves $0$. The move d3--d4 attacks
c5; every black reply except c5$\times$d4 loses at once, and after c5$\times$d4 white continues b3--b4 to
$\star$, so the reversible move d3--d4 counts as a move to $\star$. White's third move, b3--b4, loses to
c5$\times$b4. If black moves first, a7--a5 leaves $0$ and b7--b6 leaves $\star$. Black's other moves are
reversible: after a7--a6 and b7--b5 white has replies worth $2$ and $3$, both at least $\nim{2}$; and c5--c4
lets d3$\times$c4 free white's d2 pawn for $+\BIG$. Thus
$G=\game{0,\star}{0,\star}=\nim{2}$.

In position B (Fig.~\ref{fig:star2_seven_b}), if white moves first, d2--d4 leaves $0$ and c2--c3 leaves
$\star$. The moves d2--d3 and c2--c4 are reversible, since black's replies reach $-2$ and $-3$, and b4--b5
loses to a6$\times$b5. If black moves first, d5--d4 leaves $0$, and a6--a5 attacks b4: every white reply except
b4$\times$a5 loses at once, and after b4$\times$a5 black continues c6--c5 to $\star$, so a6--a5 counts as a move
to $\star$. Black's third move, c6--c5, loses to b4$\times$c5. Thus $G=\nim{2}$ again. The two analyses
match move for move because B is A with the colors reversed and the files mirrored.

\subsection{Counts, Five Files and Rings}
This subsection gives in full what Sections~\ref{sec:star2-five} and~\ref{sec:star2-rings} of the main text summarize, with every listed ring, count and check.
On ordinary boards of width at most four, in the default setting, two $\nim{2}$ positions, A and B, have seven pawns
and four have eight; one exhaustive scan of $5{,}503{,}282$ root positions gives these counts \lab{C1}, and our
engine re-checked every one \lab{C2}. With
Theorem~\ref{thm:census}, seven is the least number of pawns in a $\nim{2}$ on these boards \lab{C2}. In the
whole-board game, a scan of every position with at most seven pawns on up to nine files extends this to ordinary
boards of any width: a run of empty files acts like one, and every part has pawns of both colors and spans at most
as many files as it has pawns, so six pawns fit on eight files and seven on nine. It finds no $\nim{k}$ with $k\ge2$ with six pawns or fewer, and the only seven-pawn $\nim{2}$ are A and
B \lab{C2}.

On five files a scan of all $14{,}680{,}983$ eight-pawn positions finds $21$ of value $\nim{2}$ in the default
setting \lab{C1}, each using all five files. As whole-board games all $21$ are $\nim{2}$ by two programs \lab{C2};
as sums, $15$ keep the value. Three of the
$21$ come from the $14$ eight-pawn $\nim{2}$ rings on $\cyl{5}$ \lab{C1} by opening them onto files a--e:
\pos{W a3 b3 c2 e2 / B a6 c5 d5 e3}, \pos{W a4 b3 c2 d2 / B a7 b6 c7 e4} and
\pos{W a6 b4 c4 e3 / B a7 c7 d6 e6},
and these are $\nim{2}$ in both models. We call the first \emph{the five-file $\nim{2}$}. No scan here shows a whole-board $\nim{3}$ or higher
\lab{C1}.

By Theorem~\ref{thm:cut} these scans also cover positions with an empty file on $\cyl{3}$ to $\cyl{5}$, which
leaves rings. Two programs' scans of the rings on
$\cyl{5}$ to $\cyl{8}$ with at most eight pawns find $\nim{2}$ in the default setting only in the $14$
eight-pawn rings on $\cyl{5}$ \lab{C2}. Narrower cylinders need fewer pawns. On $\cyl{3}$ the only ring
$\nim{2}$ with six pawns in the default setting is \pos{W a2 a3 b2 / B b7 c6 c7}, and on $\cyl{4}$ the first
are \pos{W a4 b2 c2 / B b6 c5 d6 d7} and its color-reversed twin, with seven; our engine confirms both values
\lab{C2}. No smaller ring on either is a $\nim{2}$: the atlas, a C scan and an independent ring solver agree
\lab{C2}. On $\cyl{6}$ the first rings have nine pawns: a scan of every
nine-pawn ring finds six \lab{C1}, among them \pos{W a4 c4 d2 e2 / B b6 b7 d6 e5 f6}; both programs confirm
three, and the other three are their color-reversed twins \lab{C2}. On $\cyl{7}$ no ring with nine pawns is a
$\nim{2}$, and a third solver written from the rules agrees \lab{C2}; the ten-pawn rings with five pawns a side
include $63$ of value $\nim{2}$ \lab{C1}, such as \pos{W a2 a3 c3 e3 f2 / B a4 b5 d5 f5 g5}, which all three
programs confirm \lab{C2}. So the first ring $\nim{2}$ on $\cyl{n}$ has $n+3$ pawns for $n=3,\dots,7$
\lab{C2}, but on $\cyl{8}$ it has ten: no nine-pawn ring there is a $\nim{2}$ \lab{C2}, while of the ten-pawn
rings $18$ are $\nim{2}$ and two $\nim{3}$, all with five pawns a side \lab{C1}; our engine
confirms each value in both models \lab{C2}, for example \pos{W a3 b3 c3 d2 f2 / B b6 d5 e5 g4 h5}. On $\cyl{9}$
ten is again the least: no nine-pawn ring is a $\nim{2}$, but seven ten-pawn rings are \lab{C2}. From $\cyl{10}$ on, one pawn per file, the least a
ring can have, suffices: \pos{W b2 d2 e3 h3 j5 / B a4 c6 f6 g7 i7} on $\cyl{10}$ and \pos{W c5 e3 h3 i2 k2 / B a4
b7 d7 f7 g6 j6} on $\cyl{11}$ are $\nim{2}$ \lab{C2}, and with pairs of files B4, W2 inserted between files j and a
of the first, or B7, W5 between files a and b of the second, they stay $\nim{2}$ for every $n\le60$ \lab{C2}; on $\cyl{10}$ and $\cyl{11}$ no other ring with $n$ pawns is \lab{C1}.
For proofs we use \emph{walls}: runs of files with one pawn each, alternating in color, the black pawn two
ranks above the white (B4, W2: a black pawn on rank 4, a white one on rank 2); the \emph{window} is the
five or six files, core and adjacent wall files, where the reduced game is played. If a wall pawn moves, its neighbor captures it and gets a passed
pawn, so wall moves lose; with them deleted, inserting a pair of wall files leaves the \emph{reduced game}
unchanged. As a whole-board game this gives a ring $\nim{2}$ with $n+3$ pawns on
every $\cyl{n}$ with $n\ge5$, for example \pos{W b2 c3 c5 d2 f2 h2 / B a5 b5 c6 e4 g4} on $\cyl{8}$ \lab{S}: the
reduced games have $144$ and $104$ positions (odd, even $n$), and two programs check $n<17$ \lab{C2}. As a sum of components, where captures can split the ring, our engine confirms it only
to $n=12$ \lab{C1}. The same argument proves the two families with $n$ pawns above $\nim{2}$ as whole-board games for every
$n\ge10$ \lab{S}, and as sums of components they are $\nim{2}$ on $\cyl{10}$ and $\cyl{11}$ \lab{C1}. The $\cyl{3}$ ring is a $\nim{2}$ with one pawn fewer than any ordinary board of width at most four allows,
and it needs two en passant captures: without black's b4$\times$a3 it is $\uparrow$, without white's
b5$\times$c6 it is $\downarrow$, and with e.p.\ off it is $0$ \lab{C2}.

A $\nim{3}$ lies just beyond these scans, as a single component with nine pawns: B with a white a2 and a black a4 added, \pos{W a2 b4 c2 d2 / B a4 a6 a7 c6 d5}, is $\nim{3}$ on four
files, and so is A with a pair added on file e, \pos{W a4 b3 d2 d3 e4 / B a7 b7 c5 e6}, on five \lab{C2}; A beside a separate group worth
$\star$ on six files is $\nim{2}+\star=\nim{3}$ \lab{C2}. In the whole-board game no position with eight pawns or
fewer on up to eight files, ordinary or cylinder, is $\nim{3}$ \lab{C1}; as sums of components some seven-pawn positions are, but only because a $+\BIG$
cancels a $-\BIG$ (Section~\ref{sec:sums-hss}). With e.p.\ off eight pawns suffice, for example the candidate
beside a group worth $\star$ on six files, \pos{W a2 a4 c2 e2 / B a3 a6 d7 f7}, and the ring \pos{W a5 b3 c2 d2
/ B a7 b5 c6 d4} on $\cyl{4}$ under the opponent rule \lab{C2}. On $\cyl{7}$ the ten-pawn rings above include
eight of value $\nim{3}$, among them \pos{W a2 a5 c4 e2 f2 / B a7 b4 d7 e5 g7}, $\nim{3}$ in both models;
three programs agree on all eight \lab{C2}, and as sums six keep the value. On $\cyl{8}$ the ten-pawn \pos{W a4 c3 d3 e2 g2 / B b7 c6 e5 f5 h4} is $\nim{3}$, and with pairs of files B4, W2
inserted between files b and c it gives, by the same argument, a whole-board $\nim{3}$ with $n+2$ pawns on every
$\cyl{n}$ with even $n\ge8$, and a similar ring does the same for odd $n\ge11$ \lab{S}. One pawn fewer
suffices from $\cyl{11}$ on: no eleven-pawn ring there is $\nim{3}$ \lab{C1}, but the twelve-pawn
\pos{W c4 d6 e3 h3 i2 k2 / B a4 b7 d7 f7 g6 j6} is \lab{C2}, and \pos{W b2 d2 e6 f3 i3 j2 l2 / B a4 c5 e7 g7 h6
k6} on $\cyl{12}$, with pairs of files B4, W2 inserted between files l and a, and a similar ring for odd $n$,
give a $\nim{3}$ with $n+1$ pawns on every $\cyl{n}$ with $n\ge12$, by the wall argument \lab{S}. So $\nim{3}$, like $\nim{2}$, needs no more than eight ranks, but no $\nim{4}$ has turned up: none of
the $134{,}920{,}191$ ten-pawn positions on five files that form one component has a nimber above $\nim{3}$
in the default setting, nor has any of about $24$ million rings, most with at most $11$ files \lab{C1}. A ring with two
cores between walls, whenever it is a nimber, is the nim-sum of its cores \lab{C1}, so a $\nim{4}$ needs a
$\nim{4}$ core.

In the model, then, the $\nim{2}$ question has a positive answer on the ordinary $8\times8$ board: four files
and seven pawns beside T suffice, while Elkies' higher nimbers need larger boards~\cite{Elkies2002}.

\subsection{Rings on the M\"obius Strip}
The strip joins the first and last files with the ranks reversed: square $(n,r)$ is beside $(1,9-r)$. Every pawn
has a direction; white pawns start moving up and black pawns down, and a capture across the seam lands on the
mirrored rank and reverses the direction of the capturing pawn. Double step and promotion follow the direction.
A pawn is passed if no enemy pawn is strictly ahead of it, in its own direction, on its file or a neighboring
file, the seam neighbor being read in mirrored ranks; a pawn that can be taken en passant is not passed. The
whole-board game with the count rule is otherwise that of Section~\ref{sec:star2-rings}. A root ring has a pawn
on every file, no passed pawn and all pawns in their starting directions; rotating the strip does not preserve
this, so rings that differ by a rotation are counted separately.

Table~\ref{tab:mobring} gives the least number of pawns in a root ring $\nim{2}$ and one such ring for each $n$,
in ring notation: one entry per file from the first to the last, the seam lying between the last and the first.
All root rings were scanned for $n=3$ with $3$ to $10$ pawns, $n=4$ with $4$ to $9$, $n=5$ and $6$ with $n$ to
$9$, $n=7$ with $7$ and $8$, and $n=8$ with $8$ (two programs by one AI model agree on every ring). A third
program, written independently from the rules alone, agrees on the eight rings named here, on the number of
root rings, and on the number of $\nim{2}$ rings in every scan below the least, and at the least
for $n=3,4,5,8$ \lab{C2}. Two nine-pawn rings are $\nim{3}$: B67 W4 W2B6 W2B5 W5B7 on five files and
B7 W4 W23 B5 W3B7 W4B7 on six \lab{C2}.

\begin{table}[!htbp]
\centering
\caption{Least number of pawns in a root ring $\nim{2}$ on the M\"obius strip and on the cylinder, with a least ring on the strip.}
\label{tab:mobring}
\begin{tabular}{@{}cccl@{}}
\toprule
$n$ & strip & cylinder & a least ring on the strip\\
\midrule
3 & 6 & 6 & W3B6 W2B7 W4B5\\
4 & 8 & 7 & B5 W23 W2B67 W4B7\\
5 & 6 & 8 & B6 W3 B5 W2 W3B7\\
6 & 8 & 9 & B4 W2 W2 W2B6 B5 W3B7\\
7 & 8 & 10 & B6 W2 W3 B6 B7 W3 W6B7\\
8 & 8 & 10 & B7 W3 W3 B6 W3 B6 B6 W2\\
9 & 10 & 10 & B6 W2 W3 B6 B7 W3 B7 W5 W6B7\\
10 & 10 & 10 & B4 W2 B6 W2 W3 B6 B7 W3 B7 W5\\
11 & 11 & 11 & B4 W2 B6 W2 W3 B6 B7 W3 B7 W5 B7\\
\bottomrule
\end{tabular}
\end{table}

\emph{Nine to eleven files.} For $n=9$, $10$ and $11$ we scanned every root ring with $n$ pawns ($2{,}729{,}517$,
$17{,}146{,}529$ and $109{,}358{,}578$ rings) and, for $n=9$, every root ring with ten pawns ($133{,}597{,}405$ rings).
The two programs agree on every ring \lab{C1}. The independently written program agrees on every $\nim{2}$ and
$\nim{3}$ ring found, and its own scan of the nine-pawn rings on nine files also finds no $\nim{2}$ \lab{C2}. So the
least on nine files is ten pawns. Up to symmetry, ten files have two $\nim{2}$ rings with one pawn per file and
eleven files have one. On nine files one ten-pawn ring is $\nim{3}$: B5 W3 W4 B7 B7 W3 W2 W2B7 B7 \lab{C2}. For
$n=9,10,11$ the strip and the cylinder need the same number of pawns.

\emph{Families.} Three families continue these rings; $(X)^m$ stands for $m$ copies of the block $X$.
\begin{itemize}
\item $(\text{B4 W2})^m$ B6 W2 W3 B6 B7 W3 B7 W5 has $n=2m+8$ files and one pawn per file. It is $\nim{2}$ for
  $m=1,\dots,12$, that is $n=10,\dots,32$.
\item $(\text{B4 W2})^m$ B6 W2 W3 B6 B7 W3 B7 W5 B7 has $n=2m+9$ files and one pawn per file. It is $\nim{2}$ for
  $m=1,\dots,11$, that is $n=11,\dots,31$.
\item B6 W2 W3 B6 B7 W3 $(\text{B7 W5})^m$ W6B7 has $n=2m+7$ files and $n+1$ pawns. It is $\nim{2}$ for
  $m=0,\dots,12$, that is $n=7,\dots,31$.
\end{itemize}
All three families are $\nim{2}$ for every $m$ in these ranges without an upper bound: the first two for every
$m\ge1$ and the third for every $m\ge0$ \lab{S} (Section~\ref{S-star3}, Note~5: a full proof written by one AI model and checked line by line by an AI referee from another vendor). The proof is the insertion
argument of the cylinder (Section~\ref{S-ringn}): two wall files are inserted inside a wall of ten files that does not cross the seam,
and the reduced game does not change. Its base case, the reduced game at $m=5$ ($435$, $417$ and $435$ positions),
was computed by two independently written programs \lab{C2}, which also agree on the members with $m\le13$ solved directly. So one pawn per file
gives a $\nim{2}$ ring on the strip for every $n\ge10$. The members with $m=0$ of the first two
families are not $\nim{2}$, and the eight-file ring of Table~\ref{tab:mobring} is not extended by any block
insertion we tried. Since a ring has at least one pawn per file, the least number of pawns in a ring $\nim{2}$ on the
strip is $n$ for every $n\ge10$. A scan of all root rings with twelve pawns on twelve files was stopped after $50$ of
$128$ parts; the only $\nim{2}$ ring in those parts is the member of the first family.

\subsection{Boards of Other Heights}\label{S-heights}
Everything above is on eight ranks. Here the board has $R$ ranks, $5\le R\le12$: white pawns start on rank $2$,
black pawns on rank $R-1$, and the double step is allowed only from there. We solved single-component positions
without a passed pawn in which every file is occupied (rings, on a cylinder), on ordinary boards of three and four
files and cylinders of three to five files, in increasing number of pawns: about $1.8$ billion positions in all.
Table~\ref{tab:heights} gives the least number of pawns in a $\nim{2}$ and in a $\nim{3}$.

\begin{table}[!htbp]
\centering
\caption{Least number of pawns in a position of value $\nim{2}$ / $\nim{3}$ on a board of $R$ ranks. A dash: every
position of the board was solved and none has the value. ${>}p$: none with up to $p$ pawns; more pawns were not
scanned in full.}
\label{tab:heights}
\small
\begin{tabular}{@{}lcccccccc@{}}
\toprule
board & $R=5$ & $6$ & $7$ & $8$ & $9$ & $10$ & $11$ & $12$\\
\midrule
ordinary, 3 files & -- / -- & 8 / -- & 4 / -- & 10 / -- & 7 / 9 & 7 / 9 & 7 / 9 & 6 / 9\\
ordinary, 4 files & -- / -- & 6 / -- & 4 / 6 & 7 / 9 & 6 / 9 & 8 / 9 & 6 / 8 & 6 / 8\\
$\cyl{3}$ & -- / -- & -- / -- & 4 / -- & 6 / -- & 6 / ${>}12$ & 6 / 10 & 6 / 9 & 6 / 9\\
$\cyl{4}$ & -- / -- & 6 / -- & 6 / ${>}16$ & 7 / ${>}11$ & 6 / ${>}9$ & 7 / ${>}8$ & 6 / ${>}8$ & 6 / ${>}7$\\
$\cyl{5}$ & -- / -- & 8 / -- & 6 / 8 & 8 / ${>}9$ & 6 / 9 & 8 / ${>}8$ & 8 / ${>}7$ & ${>}7$ / ${>}7$\\
\bottomrule
\end{tabular}
\end{table}

No position has a nimber above $\nim{3}$ \lab{C1}. The only nimbers are $0$, $\star$, $\nim{2}$ and $\nim{3}$, with
$831$ positions of value $\nim{2}$ and $77$ of value $\nim{3}$. The least counts are not monotone in the height.
On five ranks no position is $\nim{2}$. On six ranks six pawns are enough, as in \pos{W a2 b2 d2 / B a5 c5 d5} on
four files. On seven ranks four pawns are enough, \pos{W a2 c2 / B c3 b6} on three files, and six pawns make a
$\nim{3}$, \pos{W a2 c2 d3 / B c3 b6 c6} on four files. On the cylinder of three files and nine ranks the position
with all six pawns at home, \pos{W a2 b2 c2 / B a8 b8 c8}, is $\nim{2}$. The column $R=8$ agrees with the counts of
Sections~S6.2 and~S6.3.

The solver was written for this scan. It agrees with our CGT engine on $16{,}000$ random positions, $2{,}000$ for
each height, and on each of the $48$ positions behind a least count in the table \lab{C2}. A third program,
written independently from the rules alone, repeated the complete scans for five ranks, for six ranks and for three
files on seven ranks: $2{,}381{,}595$ positions, with the same number of positions of each nimber and the same least
counts, and no nimber above $\nim{3}$ \lab{C2}. It also showed that the first solver counts $280$ impossible
positions on $\cyl{5}$ with six ranks (twenty pawns, one or two of them white); this changes no entry of the table.
For the other heights, that no smaller position exists and that no $\nim{4}$ occurs rests on the one solver
\lab{C1}. Positions with an empty file, wider boards and more pawns than the table states were not scanned.
